\documentclass[11pt]{article}
\usepackage[margin=1in]{geometry}
\usepackage{amsmath,amssymb,booktabs,hyperref}
\newcommand{\dd}{\mathrm d}
\title{Charge-even particle matching and the actual density vertex}
\author{Zimmerman programme: candidate CPB-NS-2026-10-16}
\date{16 October 2026}
\begin{document}
\maketitle
\section{The new result and its boundary}
The preceding density bridge is now connected explicitly to the original
complex carrier field by a real scalar mediator. The tree potential has
an exact positive-square representation with specified stabilizing terms.
This removes the earlier vector-current identification problem for neutral
particle--antiparticle stock at leading slow-density order. It does not
derive the mediator, its coefficients, the occupied fermion medium, a
relativistic band completion, or the galaxy law.

A separate calculation gives the density-vertex response of the neutral
gapped nodal patch. It differs from the mass/polarization-vertex kernel
computed earlier. It has zero uniform dynamical response, a computed
finite-momentum absorption, and a conservative static stability bound for
the chosen mediator contacts. Physical production, thermal history and
formation remain required. All claims here are self-reviewed and scoped;
global novelty is unestablished.

\section{A particle-level charge-even potential}
Use $\hbar=c=1$, $I=|\Phi|^2$ and the original potential
\begin{equation}
 V_c=m^2I+4m^2|X|^2+
       \left|\sqrt{\Delta_c}X+{g\over\sqrt{\Delta_c}}\Phi^2\right|^2
       +\lambda_{\rm rep}I^2,
 \qquad \Delta_c>0,\quad\lambda_{\rm rep}>0.
\end{equation}
Let $q_F$ be the signed fermion density relative to the chosen filled
bands, and let $F_F(\mathbf P,q_F)\geq0$ be the previously specified
canonical local medium energy. Introduce a real scalar $\sigma$ with
canonical positive kinetic term, mass $M_\sigma>0$, and
\begin{align}
 \delta V={}&{M_\sigma^2\over2}\sigma^2
       +a\sigma I+b\sigma q_F
       +\delta\lambda I^2+{u_0\over2}q_F^2,\\
 \delta\lambda&={a^2\over2M_\sigma^2},\qquad
 \zeta_I=-{ab\over M_\sigma^2},\qquad
 u_0={b^2\over M_\sigma^2}+u_R,\qquad
 u_R\geq{\zeta_I^2\over2\lambda_{\rm rep}}.
 \label{ns:parameters}
\end{align}
Here $[a]=E$, $[b]=1$, $[\zeta_I]=E^{-1}$ and $[u_0]=[u_R]=E^{-2}$.
The added quartic and repulsive fermion contact are chosen matching
conditions, not symmetry-protected relations or predictions of the old
action. A preferred-frame realization of the density vertex is
$b\sigma u_\mu J_F^\mu$. The medium rest frame $u^\mu$ is still an
assumption requiring a covariant physical origin. A Lorentz-scalar
$\sigma\bar\Psi\Psi$ would be a different, mass-type vertex.

There is an exact algebraic identity:
\begin{align}
 V_c+\delta V+F_F={}&m^2I+4m^2|X|^2+
 \left|\sqrt{\Delta_c}X+{g\over\sqrt{\Delta_c}}\Phi^2\right|^2\\
 &+{M_\sigma^2\over2}
       \left[\sigma+{aI+bq_F\over M_\sigma^2}\right]^2\nonumber\\
 &+\lambda_{\rm rep}
       \left[I+{\zeta_Iq_F\over2\lambda_{\rm rep}}\right]^2
 +{1\over2}\left[u_R-{\zeta_I^2\over2\lambda_{\rm rep}}\right]q_F^2
 +F_F.\nonumber
 \label{ns:positive}
\end{align}
Thus the stated density-functional extension preserves the pointwise lower
bound $V\geq m^2I+4m^2|X|^2$. The already completed Higgs-sector square
can be retained without changing this statement. The added scalar may have
nontrivial local mass shifts, but the whole potential, not an isolated
effective mass, determines this lower bound.

The dark transformation $(\Phi,X)\mapsto(e^{i\alpha}\Phi,e^{2i\alpha}X)$
still leaves all new terms invariant. With the original current convention
$j_\Phi^0=i(\Phi^*\dot\Phi-\Phi\dot\Phi^*)$, the kinetic identity is
\begin{equation}
 |\dot\Phi+im\Phi|^2
       =|\dot\Phi|^2+m^2I-mj_\Phi^0.
\end{equation}
The source has the analogous $\dot X+2imX$ identity. Therefore an extension
whose remaining Hamiltonian density is the specified nonnegative medium
energy and positive kinetic/gradient terms obeys $E\geq m|Q_c|$ relative
to its chosen vacuum. This is a conditional classical/density-functional
result. It is not a bound for an unspecified regulated fermion operator,
arbitrary finite-temperature state, gravity or the full quantum ultraviolet
theory. As before, $Q_c=0$ supplies no nonzero neutral-stock abundance or
coldness protection. No such physical closure is inferred from the square.

The exact compensating quartic in Eq.~\eqref{ns:parameters} is a convenient
scattering prescription, not a necessary equality for boundedness. More
generally set $\lambda_*=\lambda_{\rm rep}+\delta\lambda
-a^2/(2M_\sigma^2)$. The same argument works on the open parameter region
\begin{equation}
 \lambda_*>0,\qquad
 u_0-{b^2\over M_\sigma^2}>{\zeta_I^2\over2\lambda_*}.
\end{equation}
The algebraic energy condition therefore is not restricted to an exact
cancellation surface. Radiative control and actual physical matching of
that open region are still unproved.

The fermion repulsion can also have an explicit mediator origin: a healthy
massive vector coupled only to $J_F^\mu$ gives $u_0=g_F^2/M_W^2$ at
frequencies and momenta well below $M_W$. It does not require assigning a
vector charge to neutral cold mass, since the cold coupling is through
$\sigma I$. For the example below, $M_W=0.2$ GeV would give
$g_F\simeq0.1012$. This is an optional EFT realization of the chosen
repulsion, with another physical degree of freedom and mass mechanism to
count. It is not a completion of the preferred-frame band Hamiltonian or
an observationally validated dark sector.

\section{What matches to the neutral slow density}
For slowly changing, incoherent NR particle and antiparticle populations,
\begin{equation}
 \langle I\rangle_{\rm fast}\simeq {n_++n_-\over2m},
 \qquad J_c^0\simeq n_+-n_-+2(n_{X+}-n_{X-}).
\end{equation}
Unlike a vector-current source, $a\sigma I$ need not vanish for neutral
stock. Solving the soft static scalar equation gives
\begin{equation}
 \sigma\simeq-{aI+bq_F\over M_\sigma^2},\qquad
 V_{\rm eff}=V_c+\zeta_I Iq_F+{u_R\over2}q_F^2+F_F.
\end{equation}
Thus the slow charge-even bridge coefficient is
\begin{equation}
 \zeta_{\rm NR}={\zeta_I\over2m},\qquad
 K_{\rm NR,bound}={\lambda_{\rm rep}\over2m^2},\qquad
 u_R\geq{\zeta_{\rm NR}^2\over K_{\rm NR,bound}}.
\end{equation}
The $K$ expression translates the quartic lower bound using
$\langle I^2\rangle\geq\langle I\rangle^2$; it is not the actual
neutral-mixture equation of state. Occupation correlations, the opposite-charge
channel and the source population must be included to derive that EOS.
Fast coherent $2m$ terms and a mediator which cannot follow them also
prevent an unrestricted instantaneous density replacement. The fundamental
potential uses $I$ and retains those distinctions.

\section{The compensation is not an unchanged scattering theory}
The added quartic has four-point amplitude $-4\delta\lambda$. The scalar
has vertex $-ia$ and contributes two exchange trees in each elastic channel:
\begin{align}
 \delta\mathcal M_{++}&=-4\delta\lambda
       -{a^2\over t-M_\sigma^2}-{a^2\over u-M_\sigma^2},\\
 \delta\mathcal M_{+-}&=-4\delta\lambda
       -{a^2\over t-M_\sigma^2}-{a^2\over s-M_\sigma^2}.
\end{align}
For identical charges at threshold $t=u=0$, compensation is exact:
$\delta\mathcal M_{++}=0$. For a particle--antiparticle pair at threshold,
$s=4m^2,t=0$, however,
\begin{equation}
 \delta\mathcal M_{+-}={4m^2a^2\over
             M_\sigma^2(M_\sigma^2-4m^2)}.
 \label{ns:opposite}
\end{equation}
This finite annihilation-frequency remainder is precisely why soft density
matching cannot be used as all-channel vacuum matching. The formula excludes
the mediator pole; the executed family has $M_\sigma>2m$. A real scalar
also adds the physical vacuum process
\begin{equation}
 \sigma\to\Phi\Phi^*,\qquad
 \Gamma_\sigma={a^2\over16\pi M_\sigma}
                 \sqrt{1-{4m^2\over M_\sigma^2}}.
\end{equation}
The inverse process and additional scatterings must be counted if the
mediator is populated. This is an enlarged particle theory, so the old
freeze-in history and collision network cannot simply be inherited.

For the separately selected $m=2$ MeV carrier choose
$M_\sigma=20$ MeV, $a/M_\sigma=.02$, $b=-.01$ and
$u_R=\zeta_I^2/\lambda_{\rm rep}$, twice the sufficient minimum.
The tree matching gives
\begin{align}
 \delta\lambda&=2.0\,10^{-4},&
 \zeta_I&=.01\,\mathrm{GeV}^{-1},&
 \zeta_{\rm NR}&=2.5\,\mathrm{GeV}^{-2},\\
 u_R&=.00601764\,\mathrm{GeV}^{-2},&
 u_0&=.256018\,\mathrm{GeV}^{-2},&
 K_{\rm NR,bound}&=2077.23\,\mathrm{GeV}^{-2}.
\end{align}
The old same-charge threshold cross section is retained at this tree level;
the opposite-charge threshold cross section changes by about
$-3.3433\,10^{-4}$ fractionally. The vacuum scalar width is
$1.55939\,10^{-7}$ GeV and its rest lifetime is $4.22095\,10^{-18}$ s.
Neither is a medium equilibration time. Four mediator masses and three
CM speeds are tested, with exchange channels independently enumerated by
their allowed charge partitions. There is no observational fit or preferred
physical mediator mass inferred from this chosen family.

\section{Compute the density vertex, not the old mass vertex}
Take a homogeneous neutral gapped nodal patch at zero temperature.
For axial energy momentum $\mathcal Q=v k_z$ and positive imaginary
frequency $\Omega$, let $s=\sqrt{\mathcal Q^2+\Omega^2}$.
The response convention is $\delta q_F=\chi_F\delta\mu_F$ for a
Hamiltonian source $-\delta\mu_F q_F$. With nodal DOS coefficient $A$,
\begin{align}
 {\chi_F(i\Omega,\mathcal Q)\over A}
 &=\mathcal Q^2\int_0^1
       {x(1-x)\,\dd x\over\sqrt{\Delta^2+x(1-x)s^2}}\\
 &=\mathcal Q^2\left[{\Delta\over2s^2}
       +{s^2-4\Delta^2\over4s^3}
                       \arctan{s\over2\Delta}\right].
 \label{ns:density}
\end{align}
The continuous limits are understood. In particular
\begin{equation}
 \chi_F(i\Omega,0)=0,\qquad
 \chi_F(0,\mathcal Q)\simeq {A\mathcal Q^2\over6\Delta}
 \quad(\mathcal Q\ll\Delta),\qquad
 \chi_F\big|_{\Delta=0}={A\pi\mathcal Q^2\over8s}.
\end{equation}
The uniform zero follows because the total fermion number commutes with
the unperturbed band Hamiltonian. It is not the nonzero uniform mass-vertex
response from earlier calculations. An oscillating uniform $\sigma q_F$
source cannot automatically be assigned that earlier interband damping.

An independent band-projector evaluation checks Eq.~\eqref{ns:density}.
With $\mathbf p_\pm=\mathbf p\pm\boldsymbol{\mathcal Q}/2$,
$E_\pm=\sqrt{p_\pm^2+\Delta^2}$ and $S_E=E_++E_-$,
\begin{equation}
 {\chi_F\over A}={1\over2\pi}\int\dd^2p\,
 \left[1-{\mathbf p_+\cdot\mathbf p_-+\Delta^2\over E_+E_-}\right]
 {S_E\over S_E^2+\Omega^2}.
\end{equation}
The overlap, not a mass-vertex trace, fixes the normalization. Radial and
angular refinements agree with both the closed form and Feynman-parameter
integral in six selected cases. This continuum axial-patch result requires
all external energies well below the band cutoff; it is not an isotropic,
finite-band, doped or finite-temperature medium calculation.

The functional form agrees with the established gapped Dirac density
response in Ref.~\cite{ns:pyatkovskiy}, Eqs.~(A.13) and (A.16), after
removing its electromagnetic prefactors. In its two-dimensional convention
$A_{2D}=N_f/(\pi v_F^2)$ and
$\chi_F=\epsilon_0\Pi_0/(2\pi e^2)$, with $\hbar=1$ and
$\mathcal Q=v_Fk$. Our nodal coefficient $A$ has its previously derived
three-dimensional patch measure instead. This comparison authenticates
the source-dependent normalization; it does not identify dark matter with
graphene or establish novelty.

For $e^{-i\omega t}$ continue
$s=\sqrt{\mathcal Q^2-(\omega+i0)^2}$. For $\omega>0$ and
$t=\sqrt{\omega^2-\mathcal Q^2}$ the absorption is
\begin{equation}
 {\operatorname{Im}\chi_F\over A}
 ={\pi\mathcal Q^2(t^2+4\Delta^2)\over8t^3}
                 \Theta(t-2\Delta)\geq0.
\end{equation}
The sign is positive for response to a chemical source. The induced
mediator stiffness contains $-b^2\chi_F$ and hence has the passive
negative imaginary part. Finite-momentum fermion excitations can receive
energy; their actual production and heat transport are still uncomputed.

For a weak external mode $\sigma=\operatorname{Re}[\sigma_k
e^{-i\omega t+i\mathbf k\cdot\mathbf x}]$ the leading mean absorbed power
density in the repulsively dressed density channel is
\begin{equation}
 \overline{\mathcal P}_{\sigma\to F}
 ={b^2\omega\over2}{\operatorname{Im}\chi_F\over
                    |1+u_0\chi_F|^2}|\sigma_k|^2\geq0.
\end{equation}
For an approximately free narrow scalar mode of energy density
$\omega^2|\sigma_k|^2/2$, its corresponding leading energy-loss rate is
$b^2\operatorname{Im}\chi_F/(\omega|1+u_0\chi_F|^2)$.
Large dispersive self-energies require the actual mode residue instead.
Both formulas vanish at $\mathcal Q=0$ in this neutral patch. They identify
an energy receiver without supplying a temperature, equilibration time,
nonlinear occupation history or vacuum energy transfer.

\section{A conservative static mediator bound}
Retain the bare repulsion $u_0q_F^2/2$ used in the particle potential.
In the neutral density channel its random-phase response is
\begin{equation}
 \chi_{F,u_0}={\chi_F\over1+u_0\chi_F},\qquad
 0\leq\chi_{F,u_0}(0,k)\leq u_0^{-1}.
\end{equation}
For the reduced homogeneous density model with bare $I$-curvature at least
$2(\lambda_{\rm rep}+\delta\lambda)$, the soft cold susceptibility obeys
$\chi_I\leq[2(\lambda_{\rm rep}+\delta\lambda)]^{-1}$.
Additional positive cold gradients, source-square curvature or pressure
only reduce this susceptibility within that comparison. The actual neutral
particle gas has not been proved to obey that density-curvature model.
Self-gravity and negative-compressibility backgrounds are excluded.
The scalar inverse static kernel then obeys
\begin{align}
 \mathcal K_\sigma(0,k)
 &=M_\sigma^2+k^2-a^2\chi_I(0,k)
                 -b^2\chi_{F,u_0}(0,k)\\
 &\geq M_\sigma^2+k^2
       -{a^2\over2(\lambda_{\rm rep}+\delta\lambda)}-{b^2\over u_0}.
 \label{ns:kernelbound}
\end{align}
The zero-$k$ lower margin is positive when the contact inequality is strict.
Indeed the algebraic $(I,q_F,\sigma)$ Hessian with the nonnegative medium
curvature omitted has determinant
\begin{equation}
 M_\sigma^2(2\lambda_{\rm rep}u_R-\zeta_I^2)>0.
\end{equation}
Its scalar Schur complement is exactly the lower margin in
Eq.~\eqref{ns:kernelbound}. For the selected family it exceeds
$0.0116126\,M_\sigma^2$. This establishes a conservative static density-channel
test under the stated susceptibility assumptions, not all coupled
inhomogeneous modes or global dynamical stability. Doped density/mass mixing,
other spatial directions, finite temperature, counterterms and full cold
momentum response are not certified by this bound.

\section{Energy exchange and the vacuum distinction}
The canonical mediator equation and its separately defined stress obey
\begin{align}
 (\Box+M_\sigma^2)\sigma&=-aI-bq_F,\\
 T_\sigma^{\mu\nu}&=\partial^\mu\sigma\partial^\nu\sigma
  -\eta^{\mu\nu}\left[{1\over2}(\partial\sigma)^2
                         -{M_\sigma^2\over2}\sigma^2\right],\\
 \partial_\mu T_\sigma^{\mu\nu}
     &=-(aI+bq_F)\partial^\nu\sigma.
\end{align}
Thus its energy loss is reciprocal work on the other sectors when their
equations and interaction energy are included. Assigning the interaction
to one component changes the split, not total energy conservation.
The absolute metric stress of the preferred-frame medium and its rest
vector remains an additional requirement.

The canonical mediator density and isotropically averaged pressure are
\begin{equation}
 \rho_\sigma={\dot\sigma^2+|\nabla\sigma|^2+M_\sigma^2\sigma^2\over2},
 \quad p_\sigma={\dot\sigma^2\over2}
                -{|\nabla\sigma|^2\over6}
                -{M_\sigma^2\sigma^2\over2}.
\end{equation}
An isolated rapidly oscillating homogeneous massive mode has average
$p_\sigma\simeq0$, whereas a homogeneous static potential term alone has
$p_\sigma=-\rho_\sigma$. Neither limiting expression makes the
density-sourced configuration into an independently conserved vacuum fluid.
Its source, interaction stress, gradients and decay must be counted.
At $I=q_F=0$ the selected tree mediator vacuum is $\sigma=0$.
An additive constant $V_0$ remains free and supplies
$T_{V_0}^{\mu\nu}=V_0\eta^{\mu\nu}$ with pressure $-V_0$.
Nothing here fixes $V_0$ or derives the acceleration normalization from it.
Particle/medium work is not a homogeneous cold-to-vacuum transfer by itself.

\section{Remaining requirements}
This continuation supplies an explicit charge-even mediator tied to
$|\Phi|^2$, an exact stabilized potential, its tree scattering consequences,
and the correct linear neutral-medium vertex. The mediator exchanges energy
internally. Adding a vacuum constant still leaves these local equations
unchanged, and no vacuum $Q$, vacuum magnitude or $\kappa$ is derived.

The next decisive obligations are radiative matching of the chosen contact
relations, the neutral mixture's actual correlations/EOS, mediator population
and the enlarged production network, and finite-momentum transport with the
actual medium state. These must precede applying the old external-trap
density profiles to this particle candidate. A physical occupied medium,
finite stock and boundaries, full metric/lensing, native cosmology and joint
empirical validation remain necessary. The theory remains open.

\section{Evidence}
The input-hashed executions are
\texttt{neutral\_scalar\_bridge\_run\_001} (55 bounded controls) and
\texttt{scalar\_density\_response\_run\_001} (140).
The analytic square and susceptibility arguments specify the claim scope;
finite controls are not proofs of the missing physical completion.

\begin{thebibliography}{9}
\bibitem{ns:pyatkovskiy} P. K. Pyatkovskiy,
Dynamical polarization, screening, and plasmons in gapped graphene,
\href{https://arxiv.org/pdf/0808.0931v2}{arXiv:0808.0931v2},
J. Phys.: Condens. Matter 21 (2009) 025506, Eqs.~(A.13), (A.16).
\end{thebibliography}
\end{document}
